\documentclass[10pt,english]{article}

\usepackage{amsfonts,url,tikz,amsmath}

% Paper setup
\evensidemargin=0in
\oddsidemargin=0in
\textwidth=6.25in
\topmargin=-0.5in
\headheight=0.0in
\headsep=0.5in
\textheight=9.0in
\footskip=0.5in

\input{institution}

\newcommand{\Interior}[1]{\ensuremath{{#1}^{\circ}}}
\newcommand{\Closure}[1]{\ensuremath{\overline{#1}}}
\newcommand{\Complement}[1]{\ensuremath{{#1}^{c}}}

\newcommand{\Expect}{\ensuremath{\mathrm{E}}}
\newcommand{\vecnot}{\underline}
\newcommand{\RealNumbers}{\ensuremath{\mathbb{R}}}
\newcommand{\RationalNumbers}{\mathbb{Q}}
\newcommand{\ComplexNumbers}{\mathbb{C}}
\newcommand{\Real}{\mathrm{Re}}
\newcommand{\Span}{\mathrm{span}}
\newcommand{\Rank}{\mathrm{rank}}
\newcommand{\Nullity}{\mathrm{nullity}}
\newcommand{\Trace}{\mathrm{tr}}
\newcommand{\Diag}{\mathrm{diag}}
\newcommand{\dd}{\mathrm{d}}
\DeclareMathOperator*{\esssup}{ess\,sup}

% Use < , > inner product
\newcommand{\inner}[2]{{\left\langle #1 \mskip2mu , #2 \right\rangle}}
\newcommand{\tinner}[2]{{\langle #1 \mskip1mu , #2 \rangle}}

% Use < | > inner product
%\newcommand{\inner}[2]{{\left\langle #1 \mskip2mu \middle| \mskip2mu #2 \right\rangle}}
%\newcommand{\tinner}[2]{{\langle #1 \mskip1mu | \mskip1mu  #2 \rangle}}



\begin{document}

\title{ECE 586: Vector Space Methods \\ Lecture 23: Frequently Asked Questions}
\author{Henry D. Pfister \\ \univ}
\date{October 18, 2026}

\maketitle

\section{Linear and Constrained Optimization}

\paragraph{Why is an unconstrained nonconstant linear objective unbounded below?}
For $c\neq0$, choosing $x=-tc$ gives $c^T x=-t\|c\|^2\to-\infty$.
Thus the infimum is $-\infty$; there is no finite minimizing point.
If $c=0$, every point is a minimizer with value zero (slide 1).

\paragraph{Does ordinary gradient descent preserve feasibility?}
No. Subtracting the gradient can leave the feasible set.
On a nonempty closed convex feasible set, a projected step $x_{k+1}=P_{\mathcal F}(x_k-\delta_k\nabla f_0(x_k))$ preserves feasibility.
Its convergence still requires appropriate hypotheses and step sizes.
The constrained-descent picture on slide 2 should be understood as using a feasible direction or projection.

\paragraph{What does the vector inequality in a linear program mean?}
The symbol $Ax\succeq b$ means a separate scalar inequality for each row: $(Ax)_i\geq b_i$.
It does not compare vector lengths.
An LP has a linear objective and affine equality or inequality constraints, such as the forms on slide 3, ``Linear Programs.''

\paragraph{Why do positive combinations of LP constraints give lower bounds?}
For constraints $Ax\geq b$, any $\lambda\geq0$ gives $\lambda^T Ax\geq\lambda^T b$.
If $A^T\lambda=c$, this yields $c^T x\geq b^T\lambda$ for every feasible $x$.
Nonnegative coefficients preserve the inequality directions.
Slides 4--5 illustrate such certificates; any claimed coefficient combination should satisfy $A^T\lambda=c$ exactly.

\paragraph{What is the difference between feasible, optimal, and an optimal value?}
A feasible point satisfies every constraint.
An optimal point achieves the smallest objective value among feasible points.
The optimal value is the infimum over feasible points, which may fail to be achieved or may be infinite.
The notation $p^*$ on slide 6 describes a value and does not itself assert existence of an optimizer.

\section{Lagrangians and Duality}

\paragraph{Why are inequality multipliers nonnegative and equality multipliers unrestricted?}
The slides use inequalities $f_i(x)\leq0$ and equalities $h_j(x)=0$.
For feasible $x$, nonnegative $\lambda_i$ ensure $\lambda_i f_i(x)\leq0$, while $\nu_jh_j(x)=0$ for any real $\nu_j$.
Consequently, $L(x,\lambda,\nu)\leq f_0(x)$, which is the basis for weak duality.

\paragraph{Is the dual function minimized only over feasible points?}
No. It is $g(\lambda,\nu)=\inf_{x\in\mathcal D}L(x,\lambda,\nu)$, using the original domain $\mathcal D$ and putting the constraints into the Lagrangian.
For each fixed multiplier pair this gives a lower bound on every feasible objective value.
Restricting the infimum to the feasible set would change the definition (slide 10, ``Lagrangian Duality'').

\paragraph{Why does weak duality hold?}
For any feasible $x$ and any $\lambda\geq0$,
\[
g(\lambda,\nu)\leq L(x,\lambda,\nu)\leq f_0(x).
\]
Taking the primal infimum and dual supremum gives $d^*\leq p^*$.
This proof does not require the primal objective or constraints to be convex; it also does not require attainment of either optimal value.

\paragraph{Why is the dual function concave even when the original problem is not convex?}
For each fixed $x$, the Lagrangian is affine in $(\lambda,\nu)$.
The pointwise infimum of these affine functions is concave.
Concavity does not imply a unique maximizing multiplier, or even that the supremum is attained.
These distinctions qualify the wording about a ``unique max value'' on slides 10--11.

\paragraph{What is strong duality, and does it always hold?}
Strong duality means $d^*=p^*$, so the best dual lower bound equals the primal optimal value.
It is not automatic for general nonlinear problems.
For a convex problem, conditions such as Slater's condition provide standard sufficient hypotheses.
A feasible primal point and feasible dual multipliers with equal values certify optimality of both.

\paragraph{What conditions belong to the KKT system?}
For differentiable objectives and constraints, the system includes primal feasibility, $\lambda_i\geq0$, stationarity $\nabla_xL(x,\lambda,\nu)=0$, and complementary slackness $\lambda_i f_i(x)=0$.
At a local optimum, necessity requires a suitable constraint qualification, as noted on slide 8.
For a convex problem with affine equalities, satisfying the full KKT system is sufficient for global optimality.

\paragraph{Why do only active inequality constraints appear in stationarity at a KKT point?}
Complementary slackness gives $\lambda_i=0$ whenever $f_i(x)<0$.
Only active constraints with $f_i(x)=0$ can have nonzero multipliers.
Their gradients combine with the equality-constraint gradients to balance the objective gradient.
An active constraint can still have multiplier zero.

\paragraph{How does the final slide obtain the dual of $\min c^Tx$ subject to $Ax=b$, $x\geq0$?}
Using the slide's equality sign convention, write $L=c^Tx+\nu^T(b-Ax)-\lambda^Tx$ with $\lambda\geq0$.
The infimum over all $x\in\mathbb{R}^n$ is finite only if $c-A^T\nu-\lambda=0$, in which case it equals $b^T\nu$.
Eliminating $\lambda$ gives the dual $\max b^T\nu$ subject to $A^T\nu\leq c$.
When the primal LP has a finite optimal solution, the dual also has an optimal solution with the same value.

\end{document}
